%!TEX root = ../thesis.tex

\chapter{Project Development: Ideal-MHD Solver}
\label{chap:development}

\section{Development priorities after Report 1}
\label{sec:ch2-priorities}

Report~1 supplied three design decisions beyond its Euler baseline. Its
saved-output audit found no strict-IEEE CPU/GPU difference for matched HLLC
cases, so Report~2 retained matched configurations and ported HLL to CUDA. Its
variation matrix found fast-math perturbations and a separate HLLC--Rusanov
method effect, motivating recorded build semantics and an HLL/HLLD method
axis. Finally, checkpoint drift and reference comparisons showed respectively
that final states can miss temporal behaviour and that fp32--fp64 discrepancy
is not reference error. Report~2 therefore used independently executed,
predeclared stopping times and kept discrepancy and error distinct.

Changes were selected by a cost-against-exposure rule: an axis entered the
matrix if it could alter the arithmetic path without altering the algorithm,
could be varied by rebuilding or reconfiguring rather than reimplementing, and
already had a matched baseline in the harness. That admitted precision, the
build-semantics axes on host and device, the HLL/HLLD choice, the device,
resolution, observation time, CFL and requested thread count. It excluded, and
this report records as its principal coverage gaps, quadruple precision, a
second compiler, a second machine and MPI reduction order, each needing an
arithmetic backend, toolchain or parallel layer the project did not have, and
GPU HLLD, since only one kernel could be held in exact semantic correspondence
with the host. The cases follow the same rule: the Alfv\'en wave carries an
exact solution, Brio--Wu a one-dimensional Riemann problem, and Orszag--Tang
and Kelvin--Helmholtz two-dimensional wave interaction and shear evolution.
Chapter~\ref{chap:methodology} records the resulting comparison matrix.

\section{Ideal-MHD and GLM additions}
\label{sec:ch2-mhd-glm}

The MHD extension retained the finite-volume MUSCL construction of
\citet{vanleer_1979} and the Hancock predictor. In nondimensional units
with magnetic permeability set to one, the ideal-MHD equations augmented by
the generalised Lagrange multiplier (GLM) scalar $\psi$ are
\begin{align}
 \partial_t\rho+\nabla\!\cdot(\rho\mathbf v)&=0, \notag\\
 \partial_t(\rho\mathbf v)+\nabla\!\cdot
 [\rho\mathbf v\mathbf v+(p+\tfrac12|\mathbf B|^2)\mathbf I
   -\mathbf B\mathbf B]&=0, \notag\\
 \partial_t\mathbf B+\nabla\!\cdot(\mathbf v\mathbf B-\mathbf B\mathbf v)
   +\nabla\psi&=0, \qquad \nabla\!\cdot\mathbf B=0, \notag\\
 \partial_tE+\nabla\!\cdot[(E+p+\tfrac12|\mathbf B|^2)\mathbf v
   -(\mathbf v\!\cdot\!\mathbf B)\mathbf B]&=0, \notag\\
 \partial_t\psi+c_h^2\nabla\!\cdot\mathbf B&=-\psi/\tau.
 \label{eq:ch2-ideal-mhd}
\end{align}
The ideal-gas closure is $p=(\gamma-1)[E-\rho|\mathbf v|^2/2-
|\mathbf B|^2/2]$ and the conserved GLM state is
\begin{equation}
 \mathbf{U}=(\rho,\rho\mathbf{v},\mathbf{B},E,\psi)^{T},\qquad
 E=\frac{p}{\gamma-1}+\frac{1}{2}\rho|\mathbf{v}|^{2}
   +\frac{1}{2}|\mathbf{B}|^{2},
 \label{eq:ch2-mhd-state}
\end{equation}
with $\psi$ energy excluded and $\mathbf v$, $\mathbf B$ retaining three
components on one- and two-dimensional grids. The two-dimensional solver reuses
the normal flux by rotating states into the $y$-normal frame, with GLM fluxes
$F(B_n)=\psi$ and $F(\psi)=c_h^2B_n$. Damping is applied once per full
Lie-split $x$--$y$ step:
\begin{equation}
 \psi^{n+1}=\begin{cases}
 \psi^{*}\exp(-\Delta t\,c_h/c_r),&c_h>0\ \text{and}\ c_r>0,\\
 \psi^{*},&c_h\leq0\ \text{or}\ c_r\leq0.
 \end{cases}
 \label{eq:ch2-glm-damping}
\end{equation}
Here $c_h$ is the stepwise maximum directional fast-wave signal speed and $c_r$
sets $c_p^2=c_hc_r$, with $\tau=c_r/c_h$ for positive parameters. The
two-dimensional default $c_r=0.18$ is the ratio identified by
\citet{dedner_2002} as a compromise between damping and stability, and the
guard makes $c_r=0$ an undamped control. GLM transports and damps discrete
divergence error but does not enforce a solenoidal field. It was retained over
eight-wave source formulations \citep{powell_1999}, constrained transport
\citep{evans_hawley_1988,balsaraSpicer1999ct,londrilloDelZanna2004uct,
gardinerStone2005ct} and constrained hyperbolic cleaning
\citep{triccoPrice2012divb} because it fits the cell-centred, directionally
split code.

The fixed $x$-then-$y$ Lie composition is first-order as a multidimensional
split, and minmod slopes reconstruct the nine conservative components. A
non-physical reconstructed face reverts both faces to the cell average; a line
update that remains non-physical is recomputed with first-order HLL fluxes; a
post-step positivity failure terminates the run with a non-zero exit code. No
density or pressure floor is applied.

Figure~\ref{fig:ch2-implementation-delta} summarises these additions.

\begin{figure}[htbp]
  \centering
  \begin{tikzpicture}[
      box/.style={draw=fpDarkBlue, rounded corners, align=center,
                  minimum height=0.75cm, text width=3.15cm, font=\footnotesize},
      arrow/.style={->, thick, draw=fpGray}]
    \node[box, text width=4.2cm] (base) at (0,0)
      {Report~1 Euler harness\\configuration, reconstruction and output};
    \node[box] (state) at (-3.8,-2.15)
      {Nine-variable ideal-MHD state\\MHD flux and fast waves};
    \node[box] (glm) at (0,-2.15)
      {Rotated two-dimensional sweeps\\GLM transport and damping};
    \node[box] (flux) at (3.8,-2.15)
      {HLL baseline\\CPU HLLD analysis path};
    \node[box] (device) at (-2.0,-4.55)
      {CPU execution and OpenMP build control\\tested CUDA HLL mirror};
    \node[box] (gates) at (2.0,-4.55)
      {Physical-state, completion\\and regression checks};
    \draw[arrow] (base) -- (state);
    \draw[arrow] (base) -- (glm);
    \draw[arrow] (base) -- (flux);
    \draw[arrow] (state) -- (device);
    \draw[arrow] (glm) -- (device);
    \draw[arrow] (glm) -- (gates);
    \draw[arrow] (flux) -- (gates);
  \end{tikzpicture}
  \caption[Ideal-MHD implementation delta]{\textbf{Ideal-MHD implementation
  delta.} Added numerical, execution and acceptance components.}
  \label{fig:ch2-implementation-delta}
\end{figure}

\section{HLL and HLLD solver paths}
\label{sec:ch2-riemann}

HLL uses the left and right MUSCL--Hancock predicted face states. For
$k\in\{L,R\}$, its fast magnetosonic speed and Davis-type bounds
\citep{davis_1988} are
\begin{align}
 a_k^2&=\gamma p_k/\rho_k,&
 c_{A,k}^2&=|\mathbf B_k|^2/\rho_k,\notag\\
 c_{An,k}^2&=B_{n,k}^2/\rho_k,\notag\\
 c_{f,k}^2&=\tfrac12\!\left[a_k^2+c_{A,k}^2+
 \sqrt{(a_k^2+c_{A,k}^2)^2-4a_k^2c_{An,k}^2}\right],\notag\\
 S_L^{H}&=\min(v_{n,L}-c_{f,L},v_{n,R}-c_{f,R},-c_h),&
 S_R^{H}&=\max(v_{n,L}+c_{f,L},v_{n,R}+c_{f,R},c_h).
 \label{eq:ch2-hll-speeds}
\end{align}
Thus the fan contains the GLM waves $\pm c_h$. Since $c_h$ is a domain maximum
of $|v_n|+c_f$, it bounds every local candidate above, so unless
reconstruction pushes a face state past the stepwise cell-centred maximum the
clamp is never beaten and $S_L^{H}=-c_h$, $S_R^{H}=+c_h$ identically. With
$\mathbf F_k=\mathbf F(\mathbf U_k)$, the standard flux is
\begin{equation}
 \mathbf F_{H}=\begin{cases}
 \mathbf F_L,&0\le S_L^{H},\\
 \dfrac{S_R^{H}\mathbf F_L-S_L^{H}\mathbf F_R+
 S_L^{H}S_R^{H}(\mathbf U_R-\mathbf U_L)}{S_R^{H}-S_L^{H}},
   &S_L^{H}<0<S_R^{H},\\
 \mathbf F_R,&S_R^{H}\le0.
 \end{cases}
 \label{eq:ch2-hll-flux}
\end{equation}
This compact two-wave solver \citep{harten_lax_vanleer_1983} is the common
CPU/CUDA MHD path, but under the clamp its central branch collapses to
$\tfrac12(\mathbf F_L+\mathbf F_R)-\tfrac{c_h}{2}(\mathbf U_R-\mathbf U_L)$,
the Rusanov form \citep{toro2009} carrying one global wave speed instead of a
local estimate \citep{einfeldt_1988}. The tested HLL path is therefore
uniformly more dissipative than a locally estimated HLL, and the HLL/HLLD axis
separates a globally clamped two-wave flux from HLLD rather than two local
solvers; later chapters read its rates and discrepancies on that basis. The
clamp also forces $S_L^{H}<0<S_R^{H}$, so the outer branches of
Equation~\ref{eq:ch2-hll-flux} are unreachable and strict inequalities there
change nothing. HLLD has no such clamp, so its $S_L^*$, $S_R^*$ and $S_M$
ties remain reachable.

HLLD, or HLL with Discontinuities, instead uses the plain Davis outer speeds
\begin{equation}
 S_L=\min(v_{n,L}-c_{f,L},v_{n,R}-c_{f,R}),\qquad
 S_R=\max(v_{n,L}+c_{f,L},v_{n,R}+c_{f,R}),
 \label{eq:ch2-hlld-outer-speeds}
\end{equation}
without the $\pm c_h$ clamp. Following \citet{miyoshi_kusano_2005}, its
five-wave fan $S_L$, $S_L^*$, $S_M$, $S_R^*$, $S_R$ restores the Alfv\'en and
contact discontinuities and yields four intermediate fluxes. The GLM pair is
solved first:
\begin{equation}
 B_n^*=\tfrac12(B_{n,L}+B_{n,R})-\frac{\psi_R-\psi_L}{2c_h},\qquad
 \psi^*=\tfrac12(\psi_L+\psi_R)-\frac{c_h}{2}(B_{n,R}-B_{n,L}).
 \label{eq:ch2-hlld-glm-star}
\end{equation}
Every returned flux re-imposes $F(B_n)=\psi^*$ and
$F(\psi)=c_h^2B_n^*$. For $p_T=p+|\mathbf B|^2/2$,
$m_k=S_k-v_{n,k}$ and $D=m_R\rho_R-m_L\rho_L$, the key intermediate speeds are
\begin{align}
 S_M&=\frac{m_R\rho_Rv_{n,R}-m_L\rho_Lv_{n,L}-p_{T,R}+p_{T,L}}{D},&
 \rho_k^*&=\frac{\rho_km_k}{S_k-S_M},\notag\\
 S_L^*&=S_M-\frac{|B_n^*|}{\sqrt{\rho_L^*}},&
 S_R^*&=S_M+\frac{|B_n^*|}{\sqrt{\rho_R^*}}.
 \label{eq:ch2-hlld-speeds}
\end{align}
Fourteen guards test non-finite data, intermediate states and fluxes,
non-positive $\rho_k^*$, and $D$, $S_k-S_M$ and the Alfv\'en denominators
against a relative $64\varepsilon_{\mathrm{mach}}$ floor, with
$\varepsilon_{\mathrm{mach}}$ the machine epsilon of the active precision and a
minimum scale of one. Failure sends only that interface to HLL on the same
GLM-split states. Otherwise matched HLL/HLLD runs share reconstruction, updates
and cleaning, so the selected Riemann flux is the method axis.

\section{CPU, OpenMP build control, and CUDA implementation}
\label{sec:ch2-execution}

OpenMP is a shared-memory threading interface, but the MHD CPU path has no
work-sharing directives, so requests of 1, 2, 4 or 8 threads cannot alter the
decomposition. Nor could reduction order change the result under work sharing:
the only reduction here, the stepwise maximum for the CFL condition and $c_h$,
is exactly associative on finite IEEE-754 values, and order dependence would
need an accumulating reduction, absent outside diagnostics. These runs are a
null control.

CUDA, NVIDIA's GPU platform, provides an independent path mirroring CPU HLL
semantics in fp32 and fp64: boundaries, conservative minmod/Hancock
reconstruction, GLM cleaning, sweep and step order, the
Courant--Friedrichs--Lewy (CFL) time-step calculation and per-cell reversion.
CUDA obtains $c_h$ by a two-stage parallel maximum tree rather than a
sequential scan, which associativity again leaves harmless.

The CUDA MHD kernel translation unit is compiled with \texttt{--fmad=false},
scoped to that file, because MSVC's default /fp:precise does not contract
multiply--add on x64 \citep{microsoftMsvcFp2021}, while the other relevant nvcc
defaults already conform to IEEE-754 \citep{whitehead_fitflorea_2011}. Contraction
is thus the only device transformation to suppress, agreement is a property of
the declared build rather than of the hardware, and both the nvcc default and
\texttt{--use\_fast\_math} are exposed so they can be measured
(Section~\ref{sec:ch5-hardware}) instead of assumed.

CUDA omits HLLD, per-line fallback and post-step rejection.
\citet{bard_dorelli_2014} supplies context for a related GPU scheme, not a
performance baseline across different hardware and optimisation choices.

\section{Testing and acceptance criteria}
\label{sec:ch2-gates}

Unit tests isolate state conversion, pressure, wave-speed estimates, fluxes and
directional symmetry; integration tests exercise complete CPU and CUDA updates
including GLM cleaning, boundaries and transverse invariance; regression tests
preserve configuration and output contracts. A run enters measurement only if
it reaches the requested final time with a positive step count, finite states,
positive density and pressure, and fresh output and completion metadata, so
incomplete or non-physical runs are rejected before aggregation.
